%=====================================================================
%  Synthetic Data Generation: Monte Carlo Simulator
%  Methodology section. Conventions match the thesis derivation:
%  alpha^Q is the risk-neutral long-run convenience yield; the closed
%  form F_hat(tau,S,delta) is as derived and verified in the preceding
%  section.
%=====================================================================

\section{Synthetic Data Generation}
\label{sec:synthetic-data}

The inversion target of this work is the latent convenience yield
$\delta_t$, which is not observable from market data. To evaluate the
PINN's ability to recover $\delta_t$, a controlled environment is
required in which the true latent path is known by construction. This
section describes a Monte Carlo simulator that generates synthetic
futures term structures from known parameters and a stored ground-truth
$\delta_t$ path, against which the inversion can be benchmarked.

\subsection{Design rationale}

The simulator follows an \emph{Option A} architecture: physical-measure
state paths are simulated by discretising the $\mathbb{P}$-measure
stochastic differential equations, and futures prices are then obtained
from these paths using the closed-form $\mathbb{Q}$-measure pricing
formula derived in Section~\ref{sec:closed-form}, rather than by nested
Monte Carlo simulation of the risk-neutral expectation. This choice is
deliberate. The closed form returns the exact futures price conditional
on the simulated state $(S_t, \delta_t)$, so the synthetic prices carry
no Monte Carlo pricing error. Any deviation between the clean prices and
the observable data is therefore introduced \emph{deliberately} through a
measurement-noise layer (Section~\ref{sec:obs-layer}), whose magnitude is
a single controllable parameter. This separation is what makes the
simulator a useful instrument for the identifiability analysis of
Section~\ref{sec:identifiability}: the signal-to-noise ratio of the
inverse problem can be dialled directly, which a nested-simulation
approach with its own irreducible sampling noise would not permit.

\subsection{Physical-measure path simulation}

The real-world joint dynamics of the spot price and convenience yield are
the $\mathbb{P}$-measure system established in
Equations~\ref{sde1}--\ref{sde2}. The spot price is simulated in
logarithmic coordinates $X_t = \ln S_t$ to enforce positivity and to
eliminate the multiplicative-noise discretisation bias of a direct
Euler--Maruyama scheme on $S_t$. Applying It\^o's Lemma to $X = \ln S$
under $\mathbb{P}$ gives
\begin{equation}
    dX_t = \left( \mu - \delta_t - \tfrac{1}{2}\sigma_1^2 \right) dt
           + \sigma_1 \, dW_1^{\mathbb{P}},
    \label{eq:log-spot-sde}
\end{equation}
which is integrated jointly with the Ornstein--Uhlenbeck convenience
yield $d\delta_t = \kappa(\alpha^{\mathbb{P}} - \delta_t)\,dt + \sigma_2\,
dW_2^{\mathbb{P}}$ by an explicit Euler--Maruyama scheme on a uniform grid
$\{t_n = n\,\Delta t\}_{n=0}^{N}$ with $\Delta t = T/N$:
\begin{align}
    \delta_{n+1} &= \delta_n + \kappa(\alpha^{\mathbb{P}} - \delta_n)\Delta t
                    + \sigma_2 \, \Delta W_{2,n}, \\
    X_{n+1} &= X_n + \left( \mu - \delta_n - \tfrac{1}{2}\sigma_1^2 \right)\Delta t
               + \sigma_1 \, \Delta W_{1,n}.
\end{align}
The correlated Brownian increments are generated from independent
standard normals $z_{1,n}, z_{2,n} \sim \mathcal{N}(0,1)$ via the
Cholesky factor of the correlation matrix
$\left[\begin{smallmatrix} 1 & \rho \\ \rho & 1 \end{smallmatrix}\right]$,
\begin{equation}
    \Delta W_{1,n} = \sqrt{\Delta t}\, z_{1,n}, \qquad
    \Delta W_{2,n} = \sqrt{\Delta t}\left( \rho\, z_{1,n}
                     + \sqrt{1 - \rho^2}\, z_{2,n} \right),
\end{equation}
which reproduces $\mathbb{E}[\Delta W_{1,n}\Delta W_{2,n}] = \rho\,\Delta t$.
The full simulated path $\{\delta_n\}_{n=0}^{N}$ is retained as the
ground-truth latent target; it is the quantity unavailable from real data
and against which the inversion $\hat{\delta}_\theta$ is scored.

\subsection{Pricing layer}

At each simulated state $(S_n, \delta_n)$ the observable futures term
structure is computed by evaluating the closed form of
Section~\ref{sec:closed-form} across a fixed set of maturities
$\{\tau_m\}_{m=1}^{M}$:
\begin{equation}
    \ln \hat{F}(\tau_m, S_n, \delta_n)
      = \ln S_n + B(\tau_m)\,\delta_n + A(\tau_m),
\end{equation}
with $A(\tau)$ and $B(\tau)$ as derived. Because only $\alpha^{\mathbb{Q}}$
enters the pricing formula, the measure-change shift
$\alpha^{\mathbb{Q}} = \alpha^{\mathbb{P}} - \sigma_2\lambda_2/\kappa$
(Equation~\ref{sde2-Q}) is applied when passing parameters from the
simulation layer to the pricing layer; the physical $\alpha^{\mathbb{P}}$
and the market price of convenience-yield risk $\lambda_2$ govern the path
dynamics, while their risk-neutral image $\alpha^{\mathbb{Q}}$ governs the
prices.

\subsection{Observation layer}
\label{sec:obs-layer}

Real futures quotes deviate from any model price through bid--ask spreads,
non-simultaneity of observations, and data error. Following the
measurement-equation structure of the Kalman-filter formulation, these are
represented as additive, serially and cross-sectionally uncorrelated
Gaussian disturbances in log-price space:
\begin{equation}
    \ln F^{\mathrm{obs}}_{n,m}
      = \ln \hat{F}(\tau_m, S_n, \delta_n) + \varepsilon_{n,m},
    \qquad \varepsilon_{n,m} \sim \mathcal{N}(0, \sigma_\varepsilon^2).
\end{equation}
The noise scale $\sigma_\varepsilon$ is the single dial controlling the
signal-to-noise ratio of the synthetic inverse problem; setting
$\sigma_\varepsilon = 0$ recovers the exact, noise-free panel.

\subsection{Numerical verification}
\label{sec:sim-verification}

The simulator is validated before use by four checks. First, the closed
form satisfies its initial condition $\hat{F}(0, S, \delta) = S$ to
machine precision. Second, and definitively, the closed form is
substituted into the pricing PDE (Equation~\ref{GS-PDE-tau}) by automatic
differentiation at randomly sampled collocation points; the residual is of
order $10^{-14}$ in double precision, the code-level confirmation of the
analytic verification against Schwartz (1997). Third, the ensemble
statistics of the simulated $\delta_t$ paths reproduce the
Ornstein--Uhlenbeck mean $\alpha^{\mathbb{P}}$ and stationary variance
$\sigma_2^2/(2\kappa)$. Fourth, the simulated term structures exhibit the
correct qualitative regimes: backwardation at high convenience yield and
contango at low convenience yield, consistent with $B(\tau) < 0$.

\paragraph{Precision note.}
The PDE residual floor is approximately $10^{-5}$ in single precision and
falls to machine epsilon only in double precision. This is a property of
the floating-point arithmetic, not of the analytic solution, but it has a
practical consequence for the inversion: a float32 PINN cannot drive its
PDE-residual loss below this floor, which should be accounted for when
setting convergence tolerances.
